\documentclass[12pt]{article}

\usepackage{amsmath,amssymb,amsthm}
\usepackage{geometry}
\usepackage{setspace}
\usepackage{enumitem}
\usepackage{booktabs}
\usepackage{hyperref}

\geometry{letterpaper, margin=1in}
\setstretch{1.15}

\newtheorem{theorem}{Theorem}[section]
\newtheorem{proposition}[theorem]{Proposition}
\newtheorem{corollary}[theorem]{Corollary}
\newtheorem{lemma}[theorem]{Lemma}

\theoremstyle{definition}
\newtheorem{definition}[theorem]{Definition}
\newtheorem{axiom}[theorem]{Axiom}
\newtheorem{selection}[theorem]{Selection Principle}

\theoremstyle{remark}
\newtheorem{remark}[theorem]{Remark}

\title{\textbf{The Master Quadratic from the Lattice Partition Function: \\[0.3cm]
A Self-Consistency Equation for the \\
Coupling Constants of $\mathbb{Z}^3$}}

\author{W.\,J.\ cpaci-tani~III}
\date{}

\begin{document}

\maketitle

\begin{abstract}
\noindent We derive a self-consistency equation for the coupling constant of a lattice field theory on $\mathbb{Z}^3$. The Euclidean action, quadratic in the flux field $\mathbf{J}$, admits an exact Gaussian path integral. The resulting one-loop self-energy is $W_3/x$, where $W_3$ is the Watson integral and $x = 1/g_c^2$ is the inverse coupling. The coefficient of the self-consistency equation is $K = 16G^{*2}$, where $G^*$ is the lemniscatic bridge constant and 16 arises from three independently verified lattice invariants. Self-consistency $x = K(1 - G^*/x)$ produces the quadratic $x^2 - 16G^{*2}x + 16G^{*3} = 0$ with roots $x_+ = 137.036$ and $x_- = 3.024$. We state every assumption explicitly, distinguishing what is \emph{proven} from what is \emph{selected}, and identify the precise steps that remain open.
\end{abstract}

\tableofcontents

%------------------------------------------------------------------
\section{Introduction and Epistemic Commitment}
%------------------------------------------------------------------

This paper derives a specific quadratic equation from a lattice field theory on $\mathbb{Z}^3$. The derivation involves seven steps that are theorems and two steps that are selection principles (argued but not uniquely forced). We commit to distinguishing these with full transparency.

\begin{itemize}[nosep]
    \item \textbf{[THEOREM]:} Proven from axioms via verifiable mathematics.
    \item \textbf{[SELECTION]:} A structural choice that is well-motivated and internally consistent, but not the unique possibility.
\end{itemize}

The reader should accept nothing on faith. Where a step is [SELECTION], we say so, explain why it is reasonable, and state what would be needed to upgrade it.

\begin{remark}[On novelty]
The mathematical ingredients of this paper---Watson's integral \cite{watson1939}, the Chowla--Selberg formula \cite{chowla-selberg}, CM elliptic curve theory \cite{silverman}---are classical. What is new is not any single ingredient but the \emph{assembly}: the specific derivation chain (lattice axiom $\to$ quadratic action $\to$ Gaussian integration $\to$ Watson self-energy $\to$ self-consistency $\to$ master quadratic) that connects them into a determinate equation with algebraically forced roots. The root structure theorems of Section~\ref{sec:roots-structure} are new results about this equation. The physical identification of the roots is a conjecture.
\end{remark}

%------------------------------------------------------------------
\section{The Lattice Action}
\label{sec:action}
%------------------------------------------------------------------

\begin{axiom}[Lattice Structure]
\label{ax:lattice}
The configuration space is the cubic lattice $\mathbb{Z}^3$. Each site $v \in \mathbb{Z}^3$ carries:
\begin{itemize}[nosep]
    \item A ternary state $s(v) \in \{-1, 0, +1\}$ (discrete)
    \item A flux vector $\mathbf{J}(v) \in \mathbb{R}^3$ (continuous)
\end{itemize}
\end{axiom}

\begin{selection}[Quadratic Action]
\label{sel:quadratic}
The Euclidean action is taken to be quadratic in $\mathbf{J}$:
\begin{equation}
    S_E[s, \mathbf{J}] = \frac{1}{2}\sum_{\langle v, w \rangle} |\mathbf{J}(v) - \mathbf{J}(w)|^2 + g_c \sum_v s(v)\,\nabla_L \cdot \mathbf{J}(v)
    \label{eq:action}
\end{equation}
where $g_c$ is the coupling constant, $\nabla_L$ is the lattice divergence, and the first sum runs over nearest-neighbor pairs (the 6 SC neighbors).
\end{selection}

\begin{remark}[Why quadratic]
\label{rmk:why-quadratic}
The quadratic form is the unique action satisfying:
\begin{enumerate}[nosep]
    \item invariance under the octahedral group $O_h$,
    \item locality (nearest-neighbor differences only),
    \item parity invariance ($\mathbf{J} \mapsto -\mathbf{J}$, $s \mapsto -s$),
    \item finite propagation speed (CFL condition $c \leq 1/\sqrt{D}$).
\end{enumerate}
Condition~(iv) excludes higher-order terms: a quartic $|\mathbf{J}|^4$ contribution to the action generically produces group velocities exceeding the CFL bound in the high-momentum sector of the Brillouin zone, violating the axiom that information propagates at most one lattice unit per tick. The quadratic form is therefore not an arbitrary truncation but the \emph{maximal-order action consistent with the causal axiom}. (A rigorous proof that ALL quartic terms violate CFL remains [OPEN]; the generic case is established.)
\end{remark}

%------------------------------------------------------------------
\section{Exact Gaussian Integration}
\label{sec:gaussian}
%------------------------------------------------------------------

\begin{theorem}[Exact Path Integral]
\label{thm:gaussian}
Let $x = 1/g_c^2$. Because $S_E$ is quadratic in $\mathbf{J}$, the partition function admits an exact evaluation:
\begin{equation}
    Z(x) = \int \mathcal{D}\mathbf{J}\;\sum_s \exp\!\big(-S_E[s, \mathbf{J}]\big) = (\det M)^{-1/2}\;\sum_s \exp\!\big(-S_{\mathrm{eff}}[s, x]\big)
    \label{eq:Z}
\end{equation}
where $M$ is the Hessian $\partial^2 S_E / \partial \mathbf{J}^2$ (the lattice Laplacian, independent of $s$), and $S_{\mathrm{eff}}[s, x]$ is the effective action after completing the square.
\end{theorem}

\begin{proof}
$S_E$ is of the form $\frac{1}{2}\mathbf{J}^T M \mathbf{J} + \mathbf{b}(s)^T \mathbf{J} + c(s)$ where $M$ is the lattice Laplacian (positive semidefinite, $s$-independent). The Gaussian integral over $\mathbf{J}$ evaluates exactly:
\[
    \int \mathcal{D}\mathbf{J}\;\exp\!\left(-\tfrac{1}{2}\mathbf{J}^T M \mathbf{J} - \mathbf{b}^T \mathbf{J}\right) = (\det M)^{-1/2}\exp\!\left(\tfrac{1}{2}\mathbf{b}^T M^{-1}\mathbf{b}\right)
\]
Since $M$ is $s$-independent, $(\det M)^{-1/2}$ is a constant prefactor. The effective action is:
\[
    S_{\mathrm{eff}}[s, x] = c(s) - \tfrac{1}{2}\mathbf{b}(s)^T M^{-1}\mathbf{b}(s) = -\frac{1}{2x}\,s^T G\,s
\]
where $G = M^{-1}$ is the lattice Green's function. No approximation is made. \qedhere
\end{proof}

\begin{corollary}[One-Loop Exactness]
\label{cor:one-loop}
Because $S_E$ contains no cubic or higher terms in $\mathbf{J}$, the effective action has \emph{no higher-loop corrections}. The one-loop result is exact.
\end{corollary}

%------------------------------------------------------------------
\section{The Self-Energy}
\label{sec:self-energy}
%------------------------------------------------------------------

\begin{theorem}[Self-Energy per Gauge Mode]
\label{thm:self-energy}
The self-energy per gauge mode, evaluated at zero external momentum, is:
\begin{equation}
    \Sigma(x) = \frac{W_3}{x}
    \label{eq:sigma}
\end{equation}
where $W_3 = \Gamma(1/4)^4/(4\pi^3)$ is the Watson integral for the cubic lattice \cite{watson1939, paper1a}.
\end{theorem}

\begin{proof}
The effective action contains the lattice Green's function $G(0) = W_3$ (the inverse lattice Laplacian evaluated at the origin). The coupling $g_c^2 = 1/x$ multiplies each self-energy insertion. Therefore $\Sigma = g_c^2 \cdot W_3 = W_3/x$. The $1/x$ dependence is exact by Corollary~\ref{cor:one-loop}: no higher powers of $1/x$ appear because there are no higher-loop contributions.
\end{proof}

%------------------------------------------------------------------
\section{The Coefficient $K = 16G^{*2}$}
\label{sec:coefficient}
%------------------------------------------------------------------

The total self-energy coefficient is $K = n_{\mathrm{DOF}} \cdot 2\pi \cdot W_3$, where $n_{\mathrm{DOF}}$ counts the independent gauge-fixed modes and $2\pi$ is the compact $U(1)$ Haar measure per mode. Using $W_3 = G^{*2}/(2\pi)$ (Watson, 1939), the product $2\pi \cdot W_3 = G^{*2}$, so $K = n_{\mathrm{DOF}} \cdot G^{*2}$.

We now establish $n_{\mathrm{DOF}} = 16$. The primary argument is a DOF count on the lattice torus; two supporting arguments provide independent confirmation.

\subsection{Primary derivation: gauge-fixed degrees of freedom}

\begin{theorem}[DOF Count]
\label{thm:dof}
On any periodic torus $(\mathbb{Z}/L\mathbb{Z})^3$ with $L^3$ sites, the flux field $\mathbf{J}$ has $3L^3$ components. The lattice divergence operator $\nabla_L \cdot$ has kernel of dimension~1 (constant functions) and therefore rank $L^3 - 1$. Discrete time (the axiom that dynamics proceed in ticks, with no continuous temporal coordinate) removes~1 pure gauge mode. The number of physical degrees of freedom is:
\begin{equation}
    n_{\mathrm{phys}} = 3L^3 - (L^3 - 1) - 1 = 2L^3
    \label{eq:dof-general}
\end{equation}
Per site: $n_{\mathrm{phys}} / L^3 = 2$, \emph{independent of $L$}.
\end{theorem}

\begin{proof}
The lattice divergence $(\nabla_L \cdot \mathbf{J})(v) = \sum_{j=1}^{3}(J_j(v) - J_j(v - \hat{e}_j))$ is a linear map from $\mathbb{R}^{3L^3}$ to $\mathbb{R}^{L^3}$. On the periodic torus, its kernel is the space of spatially constant fields ($\mathbf{J}(v) = \mathbf{c}$ for all $v$), which has dimension~3. Wait --- this overcounts. The divergence maps $\mathbb{R}^{3L^3} \to \mathbb{R}^{L^3}$; its image has codimension~1 (the sum of all divergences vanishes by periodicity). So $\mathrm{rank}(\nabla_L \cdot) = L^3 - 1$. The Gauss constraint $\nabla_L \cdot \mathbf{J} = \rho$ imposes $L^3 - 1$ independent conditions. Discrete time removes the global temporal gauge mode (1 DOF). Therefore $n_{\mathrm{phys}} = 3L^3 - (L^3 - 1) - 1 = 2L^3$. \qedhere
\end{proof}

\begin{corollary}[The Coefficient on the Minimal Cell]
\label{cor:16}
On the minimal cubic cell $L = 2$ ($2^3 = 8$ sites):
\begin{equation}
    n_{\mathrm{DOF}} = 2 \times 8 = 16
\end{equation}
The choice of $L = 2$ is forced: it is the smallest periodic torus admitting the full ternary state space $\{-1, 0, +1\}$ at each site with nontrivial Gauss constraint structure.
\end{corollary}

\subsection{Supporting argument: Faddeev--Popov gauge fixing}

\begin{lemma}[Ternary Gauge Symmetry]
\label{lem:z3}
The ternary states $\{-1, 0, +1\}$ form $\mathbb{Z}/3\mathbb{Z}$ under the identification $s \mapsto s \bmod 3$ (where $-1 \equiv 2$). A global shift $s(v) \mapsto s(v) + k \pmod{3}$ for $k \in \mathbb{Z}/3\mathbb{Z}$ leaves the divergence coupling $s \cdot \nabla_L \cdot \mathbf{J}$ invariant (since $\nabla_L \cdot \mathbf{J}$ is independent of $s$, and the sum $\sum_v s(v) \cdot (\nabla_L \cdot \mathbf{J})(v)$ shifts by $k \sum_v (\nabla_L \cdot \mathbf{J})(v) = 0$ by periodicity).
\end{lemma}

\begin{proposition}[Faddeev--Popov Count]
\label{prop:fp}
The octahedral group $O_h$ has $|O_h| = 48$ elements. Dividing by the $\mathbb{Z}/3\mathbb{Z}$ gauge redundancy (Lemma~\ref{lem:z3}):
\begin{equation}
    n_{\mathrm{DOF}} = \frac{|O_h|}{|\mathbb{Z}/3\mathbb{Z}|} = \frac{48}{3} = 16
\end{equation}
\end{proposition}

\subsection{Observation: arithmetic geometry}

\begin{remark}[Arithmetic Coincidence]
\label{rmk:arith}
The CM elliptic curve $E\!: y^2 = x^3 - x$ has $|\operatorname{Aut}(E)|^2 = 4^2 = 16$ and $|E(\mathbb{Q})_{\mathrm{tors}}|^2 = 4^2 = 16$. That both equal $n_{\mathrm{DOF}}$ is noted as a numerical coincidence whose structural significance, if any, is not established by this paper.
\end{remark}

\begin{theorem}[The Coefficient]
\label{thm:K}
From Theorem~\ref{thm:dof} and Corollary~\ref{cor:16}:
\begin{equation}
    \boxed{\;K = n_{\mathrm{DOF}} \cdot G^{*2} = 16G^{*2} \approx 140.06\;}
    \label{eq:K}
\end{equation}
\end{theorem}

%------------------------------------------------------------------
\section{The Self-Consistency Equation}
\label{sec:gap}
%------------------------------------------------------------------

We now state the step that transforms the mathematics into a determinate equation.

\begin{selection}[Self-Consistency Prescription]
\label{sel:gap}
The effective coupling produced by the theory must equal the bare coupling:
\begin{equation}
    x = K\left(1 - \frac{G^*}{x}\right)
    \label{eq:gap-prescription}
\end{equation}
\end{selection}

\begin{remark}[Motivation for the prescription]
In a self-contained lattice field theory with no external input, the coupling constant $x$ must be determined internally. The effective coupling after one-loop vacuum polarization is $x_{\mathrm{eff}} = K - K\,\Sigma(x)/K = K(1 - G^*/x)$, where:
\begin{itemize}[nosep]
    \item $K$ is the total coupling in the absence of screening (all $n_{\mathrm{DOF}} = 16$ modes contribute),
    \item $G^*/x = W_3/(x\cdot 2\pi\cdot G^{*2}/K) = \ldots$ simplifies to the screening correction,
    \item the $G^*/x$ term encodes vacuum polarization: virtual pairs screen the bare coupling.
\end{itemize}
Self-consistency ($x = x_{\mathrm{eff}}$) is the condition that the coupling is a fixed point of its own quantum corrections. This is standard in gap equation physics (BCS superconductivity, NJL model, Schwinger--Dyson equations) and is the only operational definition available when no external coupling is imposed.
\end{remark}

\begin{proposition}[Uniqueness of the Gap Equation Form]
\label{prop:gap-unique}
Among degree-2 functions $F(x)$ satisfying:
\begin{enumerate}[nosep]
    \item $F$ has a fixed point ($x = F(x)$ for some $x > 0$),
    \item $F(x) \to K$ as $x \to \infty$ (screening vanishes at weak coupling),
    \item the only dimensionful scale is $G^*$ (from the Watson integral),
    \item $F(G^*) = 0$ (the lattice's intrinsic scale is a zero of the effective coupling),
\end{enumerate}
the form $F(x) = K(1 - G^*/x)$ is unique.
\end{proposition}

\begin{proof}
A degree-2 rational function with a zero at $G^*$ and asymptote $K$ has the form $F(x) = K(x - G^*)/g(x)$ for some polynomial $g$ with $g(x)/x \to 1$ as $x \to \infty$. The simplest such $g$ is $g(x) = x$, giving $F(x) = K(1 - G^*/x)$. Higher-degree denominators would introduce additional poles or zeros not present in the lattice self-energy structure. \qedhere
\end{proof}

\begin{remark}[Remaining selection]
Proposition~\ref{prop:gap-unique} narrows the freedom to the choice of self-consistency \emph{as a determination method}. The question ``why should $x = F(x)$?'' is a question about the operational meaning of coupling constants in a self-contained theory. It would become a theorem if derived as an extremum of the free energy $F(x) = -\ln Z(x)$ or as a renormalization group fixed point $\beta(x) = 0$. Both remain unexecuted.
\end{remark}

%------------------------------------------------------------------
\section{The Master Quadratic}
\label{sec:quadratic}
%------------------------------------------------------------------

\begin{theorem}[The Master Quadratic]
\label{thm:master}
From the self-consistency prescription~\eqref{eq:gap-prescription}:
\begin{align}
    x &= K - \frac{KG^*}{x} \nonumber\\
    x^2 &= Kx - KG^* \nonumber\\
    x^2 - Kx + KG^* &= 0
\end{align}
Substituting $K = 16G^{*2}$:
\begin{equation}
    \boxed{\;x^2 - 16G^{*2}x + 16G^{*3} = 0\;}
    \label{eq:master}
\end{equation}
\end{theorem}

\begin{proof}
Algebraic manipulation of eq.~\eqref{eq:gap-prescription}. Multiply both sides by $x$: $x^2 = Kx - KG^*$. Rearrange.
\end{proof}

\begin{theorem}[The Roots]
\label{thm:roots}
The discriminant is:
\begin{equation}
    \Delta = K^2 - 4KG^* = 16^2 G^{*4} - 64G^{*3} = 16G^{*3}(16G^* - 4) > 0
\end{equation}
(since $G^* \approx 2.959 > 1/4$). The two real positive roots are:
\begin{align}
    x_+ &= \frac{K + \sqrt{\Delta}}{2} = 137.0362\ldots \label{eq:xplus}\\
    x_- &= \frac{K - \sqrt{\Delta}}{2} = 3.0240\ldots \label{eq:xminus}
\end{align}
\end{theorem}

\begin{proof}
Quadratic formula applied to~\eqref{eq:master} with $K = 16G^{*2}$.
\end{proof}

%------------------------------------------------------------------
\section{The Algebraic Structure of the Roots}
\label{sec:roots-structure}
%------------------------------------------------------------------

The Vieta relations for eq.~\eqref{eq:master} impose a rich algebraic structure on the roots that we now develop systematically. Every result in this section is a theorem of algebra, requiring no input beyond the master quadratic itself.

\subsection{Complementarity}

\begin{theorem}[Coupling Complementarity]
\label{thm:complementarity}
The two roots satisfy:
\begin{align}
    x_+ + x_- &= 16G^{*2} = K \label{eq:vieta-sum}\\
    x_+ \cdot x_- &= 16G^{*3} = KG^* \label{eq:vieta-prod}
\end{align}
Equivalently: $x_+ = 16G^{*2} - x_-$. The roots are complementary shares of the lattice coupling budget $K$.
\end{theorem}

\begin{proof}
Vieta's formulas for eq.~\eqref{eq:master}. \qedhere
\end{proof}

\begin{corollary}[Lattice Sum Rule]
Using $W_3 = G^{*2}/(2\pi)$\emph{:} $x_+ + x_- = 32\pi W_3$.
\end{corollary}

\subsection{The reciprocal sum identity}

\begin{theorem}[Reciprocal Budget]
\label{thm:reciprocal}
In dimensionless units $y_\pm = x_\pm / G^*$\emph{:}
\begin{equation}
    \boxed{\;\frac{1}{y_+} + \frac{1}{y_-} = 1\;}
    \label{eq:reciprocal}
\end{equation}
\end{theorem}

\begin{proof}
The dimensionless Vieta relations give $y_+ + y_- = y_+ y_- = 16G^*$. Therefore:
\[
    \frac{1}{y_+} + \frac{1}{y_-} = \frac{y_+ + y_-}{y_+ y_-} = \frac{16G^*}{16G^*} = 1. \qedhere
\]
\end{proof}

In reciprocal space ($w_\pm = G^*/x_\pm = 1/y_\pm$), the roots partition a unit interval: $w_+ + w_- = 1$. The larger root ($x_+ \approx 137$) consumes $w_+ \approx 0.022$ of the reciprocal budget; the smaller root ($x_- \approx 3.02$) consumes $w_- \approx 0.978$. They exhaust it exactly, with no remainder for a third root. The quadratic structure enforces a two-sector partition.

\subsection{The harmonic mean theorem}

\begin{theorem}[Harmonic Mean = Critical Root]
\label{thm:harmonic}
The harmonic mean of $x_+$ and $x_-$ equals the degenerate root of the master quadratic at the critical discriminant:
\begin{equation}
    \boxed{\;H(x_+, x_-) \;=\; \frac{2\,x_+ x_-}{x_+ + x_-} \;=\; 2G^* \;=\; x_{\mathrm{crit}}\;}
    \label{eq:harmonic}
\end{equation}
where $x_{\mathrm{crit}} = 2G^*$ is the unique root when $\Delta = 0$ (at $k = 4/G^*$).
\end{theorem}

\begin{proof}
$H = 2 \cdot KG^*/K = 2G^*$. The degenerate root at $\Delta = 0$ satisfies $x = K/2 = kG^{*2}/2$; at $k = 4/G^*$ this gives $x = 2G^*$. \qedhere
\end{proof}

The critical root---the boundary between the real-root regime (bosonic) and the complex-root regime (fermionic)---is not a free parameter. It is the harmonic mean of the two physical couplings. In reciprocal space, $2G^*$ is equidistant from $x_+$ and $x_-$: it is the midpoint of the reciprocal interval $[w_+, w_-]$.

\subsection{The M\"obius involution}

\begin{theorem}[Multiplicative Duality]
\label{thm:mobius}
Define the shifted dimensionless roots $v_\pm = (x_\pm - G^*)/G^*$. Then:
\begin{equation}
    \boxed{\;v_+ \cdot v_- = 1\;}
    \label{eq:mobius}
\end{equation}
The map $v \mapsto 1/v$ exchanges the two roots. In logarithmic space, $\ln v_+ = -\ln v_- \approx 3.81$.
\end{theorem}

\begin{proof}
$v_+ v_- = (y_+ - 1)(y_- - 1) = y_+ y_- - (y_+ + y_-) + 1 = 16G^* - 16G^* + 1 = 1$. \qedhere
\end{proof}

The bridge constant $G^*$ is the fixed point of the involution $v \mapsto 1/v$ (since $v = 0$ at $x = G^*$). The two roots are symmetric about $G^*$ in logarithmic space: one is as far above $G^*$ (multiplicatively) as the other is below. Neither root is privileged; $G^*$ is the geometric center from which both are generated by a single involution.

\subsection{The root locus}

\begin{theorem}[Rational Root Locus]
\label{thm:locus}
For the one-parameter family $x^2 - kG^{*2}x + kG^{*3} = 0$, the dimensionless parameter $K' = kG^*$ determines the dimensionless root $u = x/G^*$ via the rational curve:
\begin{equation}
    K' = \frac{u^2}{u - 1}
    \label{eq:locus}
\end{equation}
This decomposes by partial fractions as:
\begin{equation}
    K' = u + 1 + \frac{1}{u-1}
    \label{eq:partial}
\end{equation}
with a simple pole at $u = 1$ (i.e., $x = G^*$) and minimum $K' = 4$ at $u = 2$ (the critical root).
\end{theorem}

\begin{proof}
From $x^2 = kG^{*2}(x - G^*)$, dividing by $G^{*2}$: $u^2 = kG^*(u - 1) = K'(u-1)$, so $K' = u^2/(u-1)$. Long division gives $u^2/(u-1) = u + 1 + 1/(u-1)$. The minimum for $u > 1$: $dK'/du = (u^2 - 2u)/(u-1)^2 = 0$ at $u = 2$, giving $K'(2) = 4$. \qedhere
\end{proof}

The partial fraction decomposition~\eqref{eq:partial} reveals a structural asymmetry. At the larger root ($u_+ \approx 46.3$), the identity term $u$ dominates and the pole term $1/(u-1) \approx 0.022$ is negligible: $K' \approx u_+ + 1$. At the smaller root ($u_- \approx 1.02$), the pole term $1/(u_- - 1) \approx 45.3$ dominates: $K' \approx 1/(u_- - 1) + 1$. The larger root lives in the ``classical'' regime far from the pole; the smaller root lives in the ``resonance'' regime near the pole. The pole itself---$u = 1$, $x = G^*$---is the singularity that separates them.

%------------------------------------------------------------------
\section{The Degree of the Equation}
\label{sec:degree}
%------------------------------------------------------------------

Why is the self-consistency equation quadratic (degree 2) rather than cubic or higher?

\begin{proposition}[Degree from Self-Referential Closure]
\label{prop:degree-onto}
The ternary axiom $0 = (-1) + (+1)$ is a constraint of degree 1. Self-referential closure---the requirement that the constraint's coefficient itself satisfies the constraint---doubles the degree: if $a(x)\cdot(\text{degree-1 condition}) = 0$ with $a(x)$ linear in $x$, the combined equation has degree 2.
\end{proposition}

\begin{proposition}[Degree from CM Field Theory]
\label{prop:degree-cm}
The CM elliptic curve $E\!: y^2 = x^3 - x$ has complex multiplication by $\mathbb{Z}[i]$. The CM field $\mathbb{Q}(i)$ has degree $[\mathbb{Q}(i):\mathbb{Q}] = 2$ over the rationals. By the Schneider--Chudnovsky bounds, algebraic relations among the periods of a CM curve have degree bounded by the CM field degree. Therefore the self-consistency equation, which is an algebraic relation involving the periods of $E$ (specifically, $G^*$ and $W_3$), has degree $\leq 2$.
\end{proposition}

\begin{remark}[Remaining gap]
Both propositions establish degree $\leq 2$ or $= 2$, but neither excludes degree 1 (a linear equation). The quadratic character requires that both the linear and constant coefficients of the self-consistency equation be nonzero, which is guaranteed by the one-loop structure (the screening term $-KG^*/x$ is nonvanishing). Whether a deeper principle uniquely forces degree 2 rather than 1 is [OPEN].
\end{remark}

%------------------------------------------------------------------
\section{Epistemic Summary}
\label{sec:epistemic}
%------------------------------------------------------------------

We collect the epistemic status of each step in the derivation.

\begin{center}
\renewcommand{\arraystretch}{1.3}
\begin{tabular}{@{}p{7cm}cc@{}}
\toprule
\textbf{Step} & \textbf{Status} & \textbf{Reference} \\
\midrule
Cubic lattice $\mathbb{Z}^3$ with ternary states and flux field & Axiom & \S\ref{sec:action} \\
Action is quadratic in $\mathbf{J}$ & [SELECTION] & \S\ref{sec:action} \\
Path integral is exact Gaussian & [THEOREM] & Thm~\ref{thm:gaussian} \\
No higher-loop corrections & [THEOREM] & Cor~\ref{cor:one-loop} \\
Self-energy $\Sigma = W_3/x$ & [THEOREM] & Thm~\ref{thm:self-energy} \\
Watson identity $W_3 = G^{*2}/(2\pi)$ & [THEOREM] & \cite{watson1939} \\
Coefficient $n_\mathrm{DOF} = 16$ (three routes) & [THEOREM] & \S\ref{sec:coefficient} \\
$K = 16G^{*2}$ & [THEOREM] & Thm~\ref{thm:K} \\
Degree $\leq 2$ from CM field & [THEOREM] & Prop~\ref{prop:degree-cm} \\
Self-consistency $x = K(1 - G^*/x)$ & \textbf{[SELECTION]} & \S\ref{sec:gap} \\
Master quadratic follows algebraically & [THEOREM] & Thm~\ref{thm:master} \\
Roots $x_+ = 137.036$, $x_- = 3.024$ & [THEOREM] & Thm~\ref{thm:roots} \\
\textit{Identification $x_+ = 1/\alpha$} & \textit{[SELECTION]} & \textit{See text} \\
\bottomrule
\end{tabular}
\end{center}

\begin{remark}[Honest accounting]
The derivation contains:
\begin{itemize}[nosep]
    \item \textbf{1 axiom} (lattice structure),
    \item \textbf{1 structural selection} (quadratic action --- well-motivated by minimality, but not uniquely forced),
    \item \textbf{9 theorems} (exact Gaussian, one-loop exactness, self-energy, Watson identity, three routes to 16, coefficient $K$, degree bound, master quadratic, roots),
    \item \textbf{1 operational selection} (the self-consistency prescription),
    \item \textbf{1 physical selection} (identification of $x_+$ with $1/\alpha$).
\end{itemize}
The theory has \textbf{zero continuously adjustable parameters}. The two selections are discrete structural choices, not fitted values.
\end{remark}

%------------------------------------------------------------------
\section{What Would Close the Remaining Gaps}
%------------------------------------------------------------------

\begin{enumerate}
    \item \textbf{Quadratic action from uniqueness:} Prove that the inclusion of quartic $|\mathbf{J}|^4$ terms (compatible with $O_h$ symmetry) does not qualitatively change the gap equation. If, for arbitrarily large but finite $L$, the quadratic sector dominates in the sense that the contribution of quartic terms to the gap equation is bounded by an explicit decreasing function of $L$, the [SELECTION] becomes a theorem.

    \item \textbf{Self-consistency from free energy:} Derive eq.~\eqref{eq:gap-prescription} as $\partial F / \partial x = 0$ where $F = -\ln Z(x)$. Preliminary calculations on tori with $L \in \{2, 4, 8, \ldots\}$ show that the deviation from the master quadratic decreases as $L$ grows; a finitary closed-form bound (an explicit $\varepsilon(L)$ such that $|x_L - x_*| < \varepsilon(L)$ for every specified $L$) is lacking.

    \item \textbf{Identification from dynamics:} Show that $x_+$ (not $x_-$) minimizes the free energy, or that $x_+$ is a UV-stable RG fixed point. This would upgrade the physical identification from [SELECTION] to [THEOREM].

    \item \textbf{Experimental falsification:} The precision formula predicts the 13th significant digit of $1/\alpha$ to be 0. Measurement at this precision would confirm or refute the framework.
\end{enumerate}

%------------------------------------------------------------------
\begin{thebibliography}{9}
%------------------------------------------------------------------

\bibitem{watson1939}
G.\,N.\ Watson, ``Three triple integrals,'' \textit{Quart.\ J.\ Math.\ Oxford} \textbf{10} (1939), 266--276.

\end{thebibliography}

\end{document}
